% Zdroj: PDF priklady-leto-2025.pdf, fyzická str. 1-3. Značka: společné okruhy. Zařazeno: I4.
\section{Jednoduchý správce paměti}
\szzotazka{léto 2025}{2}{Jednoduchý správce paměti (heap alokátor, 32-bit)}

\noindent\textbf{Zadání.}\par
Předpokládejte v~celé otázce pouze kontext 32-bitového \textbf{little endian} procesoru,
\textbf{bitem 0 rozumíme LSb}.

Naším cílem je v~souvislém bloku paměti implementovat jednoduchý heap pro dynamickou
alokaci podbloků paměti (neplést s~haldou jako datovou strukturou). Celý náš heap bude mít
pevnou velikost danou při inicializaci a~nebude používat žádné optimalizace nad rámec zde
popsaného. Paměť našeho heapu si chceme organizovat jako souvislou posloupnost podbloků,
kde každý podblok má být vždy \textit{zarovnaný na 2 byty}. Každý podblok heapu je buď
\textit{volný} nebo \textit{obsazený} (naalokovaný).

\medskip
\noindent\textbf{Volné podbloky} mají následující strukturu:\par
\texttt{+0}: Size [16-bit] = velikost podbloku \textbf{včetně této hlavičky}. Jelikož
z~výše uvedeného vyplývá, že bit 0 této položky by byl vždy 0, tak tuto nulu neukládáme.
Místo toho pro nás bude mít bit 0 položky Size speciální význam -- bude označovat, zda je
podblok volný nebo obsazený. $0 = $ \textit{volný}, $1 = $ \textit{obsazený}. Tedy pro
volný podblok:\par
\texttt{bit 0: 0} $=$ volný\par
\texttt{+2}: Next [16-bit] = offset dalšího volného podbloku od začátku celého heapu --
tedy volné podbloky tvoří vlastně jednosměrně vázaný seznam. Poslední volný podblok má
v~této položce vyplněnou hodnotu \texttt{-1}.\par
\texttt{+4}: zbytek podbloku je nevyužitý a~je vyplněn nulovými byty.

\medskip
\noindent\textbf{Obsazené podbloky} mají následující strukturu:\par
\texttt{+0}: Size [16-bit] = má stejný význam jako u~volného podbloku, s~tím rozdílem, že zde:\par
\texttt{bit 0: 1} $=$ obsazený\par
\texttt{+2}: Payload -- samotná data podbloku (náš alokátor vrací všechny tyto byty
vynulované).\par
\textit{Pozor:} obsazené podbloky tedy nemají položku Next -- místo v~paměti, které
u~volného podbloku náleželo položce Next, se alokací podbloku stane součástí Payloadu
a~musí být alokátorem na haldě vynulováno.

\medskip
Pro celý heap si pamatujeme offset prvního volného podbloku od začátku celého heapu. Na
počátku se heap skládá právě z~1 volného podbloku, jehož velikost odpovídá velikosti
celého heapu. Při požadavku na alokaci nového podbloku dostaneme požadovanou velikost
payloadu v~bytech a~použijeme první dostatečně velký volný podblok na nejnižším offsetu od
začátku heapu (tj. strategie first-fit).

Předpokládejte následující kostru třídy v~C\#, kterou budete doplňovat (požadovanou
implementaci můžete napsat v~jazyce C\# nebo Java nebo C++ -- pro Javu nebo C++
předpokládejte ekvivalentní implementaci kostry v~daném jazyce s~tím, že parametry
konstruktoru a~metod si můžete vhodně upravit dle zvyklostí v~daném jazyce) (\texttt{ushort}
je 16-bitové bezznaménkové celé číslo):

\begin{lstlisting}[language={[Sharp]C}]
class Heap {
    private ushort _firstFreeOffset = 0;
    private byte[] _heap;

    public Heap(byte[] availableMemory) {
        _heap = availableMemory;
        // TODO
    }
    private ushort FindFirstFree(ushort payloadSize) { // TODO }
    private void Mark(ushort offset, bool isFree) { // TODO }

    public void SetUshort(ushort offset, ushort value) { ... }
    public ushort GetUshort(ushort offset) { ... }
    public ushort Alloc(ushort payloadSize) { ... }
    public void Free(ushort payloadOffset) { ... }
}
\end{lstlisting}

Metody \texttt{SetUshort} a~\texttt{GetUshort} \textbf{ne}programujte, ale naopak je
vhodně využijte ve svém kódu. Metoda \texttt{SetUshort} na \texttt{offset}u od začátku
heapu uloží \texttt{value} v~little endian pořadí. Metoda \texttt{GetUshort} vrátí
16-bitovou hodnotu uloženou na heapu od daného \texttt{offset}u (byty na heapu chápe
v~little endian pořadí).

Metody \texttt{Alloc} a~\texttt{Free} \textbf{ne}programujte -- tyto metody bude programovat
někdo jiný s~využitím vašich metod. Pro úkol 4 níže se hodí vědět: metoda \texttt{Alloc}
naalokuje na heapu místo pro \texttt{payloadSize} velký payload a~vrátí offset pro payload;
metoda \texttt{Free} bere offset payloadu vráceného nějakým voláním metody \texttt{Alloc}
a~podblok, do kterého payload náleží, uvolní.

\medskip
\noindent\textbf{Úkoly:}
\begin{enumerate}
\item Doplňte implementaci konstruktoru, který inicializuje prázdný heap.
\item Napište implementaci metody \texttt{FindFirstFree}, která vrátí offset prvního
  volného podbloku vhodné velikosti. Stav podbloku tato funkce nijak nemění.
\item Napište implementaci metody \texttt{Mark}, která očekává, že na \texttt{offset}u je
  platný podblok haldy, a~změní jeho stav na \uv{obsazený} nebo \uv{volný} dle parametru
  \texttt{isFree}. Funkce nebude manipulovat s~položkou Next.
\item Předpokládejte, že jsme pro náš heap dostali k~dispozici 22 bytů -- tato paměť je
  předem vyplněna nulovými byty. Předpokládejte, že provedeme s~heapem následující operace:

\begin{lstlisting}[language={[Sharp]C}]
A = Alloc(2)
B = Alloc(4)
C = Alloc(2)
Free(B)
\end{lstlisting}

Napište hexdump 22 bytů paměti konečného stavu heapu po provedení všech těchto operací --
řešení vyplňte do jedné z~níže uvedených tabulek (je jich zde redundantně více, kdybyste
své řešení potřebovali opravit).

\begin{center}
\begin{tabular}{|c|c|c|c|c|c|c|c|c|}
\hline
\textbf{offset} & +0 & +1 & +2 & +3 & +4 & +5 & +6 & +7 \\
\hline
0x0000 & & & & & & & & \\
\hline
0x0008 & & & & & & & & \\
\hline
0x0010 & & & & & & & & \\
\hline
\end{tabular}
\end{center}
\end{enumerate}

\medskip
\noindent\textbf{Náčrt řešení.}\par
Konečný stav heapu (hexdump 22 bytů):
\begin{center}
\begin{tabular}{|c|c|c|c|c|c|c|c|c|}
\hline
\textbf{offset} & +0 & +1 & +2 & +3 & +4 & +5 & +6 & +7 \\
\hline
0x0000 & 05 & 00 & 00 & 00 & 06 & 00 & 0E & 00 \\
\hline
0x0008 & 00 & 00 & 05 & 00 & 00 & 00 & 08 & 00 \\
\hline
0x0010 & FF & FF & 00 & 00 & 00 & 00 & -- & -- \\
\hline
\end{tabular}
\end{center}

\medskip
\noindent\textit{Doplňující vysvětlení (nad rámec oficiálního náčrtu):} Oficiální náčrt uvádí jen
výsledný hexdump (úkol~4). Úkoly 1--3 v~C\#: potřebná velikost podbloku je payload $+\,2$ (hlavička),
zaokrouhleno nahoru na sudé číslo.
\begin{lstlisting}[language={[Sharp]C}]
public Heap(byte[] availableMemory) {
    _heap = availableMemory;
    _firstFreeOffset = 0;
    SetUshort(0, (ushort)_heap.Length);   // Size = cely heap, bit 0 = 0 (volny)
    SetUshort(2, 0xFFFF);                 // Next = -1 (posledni volny)
    for (int i = 4; i < _heap.Length; i++) _heap[i] = 0;   // zbytek nulovy
}
private ushort FindFirstFree(ushort payloadSize) {
    ushort need = (ushort)(payloadSize + 2);
    if ((need & 1) != 0) need++;                        // zarovnani na 2 B
    ushort o = _firstFreeOffset;
    while (o != 0xFFFF) {
        ushort size = (ushort)(GetUshort(o) & 0xFFFE);  // bit 0 je priznak
        if (size >= need) return o;
        o = GetUshort((ushort)(o + 2));                 // Next
    }
    return 0xFFFF;                                      // nic vhodneho
}
private void Mark(ushort offset, bool isFree) {
    ushort size = GetUshort(offset);
    SetUshort(offset, isFree ? (ushort)(size & 0xFFFE) : (ushort)(size | 1));
}
\end{lstlisting}
Odvození hexdumpu (22~B, na počátku jeden volný podblok \texttt{@0} velikosti~22):
\texttt{Alloc(2)} potřebuje $2{+}2=4$~B, vezme \texttt{@0} a~zbytek \texttt{@4} (18~B) zůstane volný;
\texttt{Alloc(4)} potřebuje 6~B, vezme \texttt{@4}, volný zbytek \texttt{@10} (12~B); \texttt{Alloc(2)}
vezme \texttt{@10} (4~B), volný zbytek \texttt{@14} (8~B). \texttt{Free(B)} dostane offset payloadu~6,
podblok je tedy \texttt{@4}; stane se volným a~zařadí se do seznamu před \texttt{@14}. Výsledek:
A~\texttt{@0} se \texttt{Size}~$= 4\,|\,1 = 5$, volný \texttt{@4} se \texttt{Size}~6 a~\texttt{Next}~$=14=\texttt{0x0E}$,
C~\texttt{@10} se \texttt{Size}~5, volný \texttt{@14} se \texttt{Size}~8 a~\texttt{Next}~$=\texttt{0xFFFF}$.
Všechny 16-bitové hodnoty jsou little-endian (\texttt{05 00}, \texttt{0E 00}, \dots), payloady i~zbytky
volných podbloků jsou nulové.
